\clearpage\nobalance
\section*{Author Notes}
\subsection*{Acknowledgements}
I thank Prof. Luyao Zhang, Program Chair, and Prof. Ken Rogerson, Program Discussant, of the \emph{2056 Nobel and Turing Laureate Program Launch Workshop}, held September 7, 2026 at Duke Kunshan University. I also thank Jiayang Sun and Zhenning Wang, my teammates in Week 3 session A, and the COMSCI/ECON 206 class for discussion and feedback. I am responsible for the proposal and remaining errors.
\subsection*{Individual contribution}
I developed the research question, interdisciplinary mechanism, proposal, computational design, and learning-tool concept. Any shared material or feedback is identified where used.

\bibliographystyle{ACM-Reference-Format}
\bibliography{references}

\appendix
\section{Technical details and reproducibility}
\textbf{Links and version.} Hugging Face: \DemoLink. GitHub: \GitHubLink. Colab: \ColabLink. Author: Xuantong Fu. Tested: [date], version [version/commit].

\textbf{Mechanism.} Buyers face scarce goods under two rules: refundable/non-transferable versus transferable/non-refundable. Agents differ in income uncertainty and speculative participation. Reputation signals vary in reliability and visibility. Outcomes are purchase probability, surplus, speculation, and access by buyer type. The simulation resets from a documented parameter file and uses fixed seeds where stochasticity is introduced.

\textbf{Compact rule.} For each treatment: initialize heterogeneous buyers and scarce supply $\rightarrow$ reveal rule and reputation information $\rightarrow$ buyers choose/submit bids or purchase decisions $\rightarrow$ allocate goods $\rightarrow$ record purchase, surplus, speculation, and access $\rightarrow$ repeat across demand/scarcity conditions.

\textbf{Checks.} Boundary: zero speculative buyers and minimum demand. Typical: baseline demand, scarcity, and mixed buyer types. Invalid input: negative demand or an income probability outside $[0,1]$; the implementation rejects or normalizes invalid values rather than silently treating them as valid. Results from the Week 2 browser prototype are reported separately from future simulation results.

\textbf{Responsible boundary.} The learning demo uses a static browser implementation only: no Hugging Face CPU/GPU runtime, backend, database, secret, API call, or paid service. Checks include desktop/mobile layout, keyboard interaction, readable labels, reset behavior, and invalid-input handling. Privacy is not collected or stored. Population-level behavioral claims remain unverified until larger experiments are conducted.

\section{AI-assistance statement}
I used ChatGPT (GPT-5.6 Luna) during September 2026 for bounded support with proposal structure, wording, computational-design discussion, and debugging of the browser prototype. I retained responsibility for the research question, disciplinary integration, hypotheses, treatment conditions, outcome measures, and final decisions. AI suggestions were treated as drafts: alternatives that were too broad, unsupported, or inconsistent with the intended mechanism were rejected or revised. I checked code behavior by running the game and checking transitions, locked choices, reset behavior, boundary/invalid inputs, and displayed calculations. I checked academic claims against the original papers and kept synthetic or pilot outputs separate from planned findings. AI output was not used as scholarly evidence.

\section{Cumulative Intellectual Development}
This proposal develops my Intellectual Statement II through the \emph{2056 Nobel and Turing Laureate Program Launch Workshop}, September 7, 2026, Duke Kunshan University, IB 2050. Program Chair: Prof. Luyao Zhang. Program Discussant: Prof. Ken Rogerson. My session: A. Rethinking Reputation in Online Marketplaces. My role: Behavioral Scientist. My teammates: Jiayang Sun, Zhenning Wang.

The earlier project focused mainly on how reputation information changes trust and transaction decisions. Peer discussion and field evidence motivated a broader question: trust is valuable, but platform rules and financial constraints can determine who actually receives scarce opportunities. I therefore added income uncertainty, speculation, and ticketing rules, changing the project from a reputation-only problem into a joint trust-and-access problem. [Add the exact dated peer comment and its source before submission.]

\section{Field and open-source inspiration}
The \emph{Joint Interdisciplinary Field Trip---Technology, Computation, Culture \& Society in Action} took place September 4, 2026. The field evidence includes the Tencent Shanghai Office green-cloud-computing exhibit and the Shanghai Science and Technology Museum four-axis trainer. The photographs are retained with their captions and alt text from Week 2. The Tencent exhibit directly showed energy-efficiency and renewable-energy measures; my interpretation is that computational systems can change resource allocation at organizational scale. The museum exhibit directly showed an interactive aerospace-motion experience; my interpretation is that information and technology are partly effective through how people experience and understand them.

The Week 2 open-source tools were Microsoft Recommenders and the CDC Measles Outbreak Simulator. Microsoft Recommenders motivated the partner-ranking/opportunity-concentration perspective; its recommendation focus limits direct transfer to the marketplace mechanism. The CDC simulator motivated interactive comparison of intervention conditions; its disease-specific assumptions limit direct transfer to market behavior. Official links, repositories, licenses, versions, and access dates are retained in the Week 2 supporting materials.

\section{Review response and revision record}
\textbf{Strength retained:} the proposal keeps a concrete platform setting and a clear interdisciplinary mechanism.\\
\textbf{Constructive concern addressed:} [insert actual peer criticism] led me to [insert actual revision].\\
\textbf{Question clarified:} [insert actual peer question] led me to specify [mechanism, variable, or test].\\
Future peer-review results will be added without replacing earlier versions.

\clearpage\onecolumn
\section{Structured abstract and cumulative proposal map}\label{app:abstract}
\begin{table}[H]
\caption{Cumulative proposal map.}
\label{tab:abstract-map}
\renewcommand{\arraystretch}{1.18}
\begin{tabularx}{\textwidth}{@{}p{0.27\textwidth}Y@{}}
\toprule
\textbf{Proposal element} & \textbf{Current statement / change} \\
\midrule
Background & Digital markets reduce transaction frictions, but reputation, rules, and financial constraints can jointly shape trust and access. \\
Motivation and application & Scarce goods create a joint problem of information design, speculation, platform rules, and unequal purchasing capacity. \\
Three connected questions & Economics studies income constraints; computation compares rule--signal combinations; behavior studies trust and action under uncertainty. \\
Method and model assumptions & Heterogeneous-agent simulation with two ticketing rules, variable demand/scarcity/speculation, and reputation conditions. \\
Results or expected results & The Week 2 game provides preliminary observed outputs; PS1 simulation results remain planned until reproduced in the notebook. \\
Intellectual merits & The synthesis asks when information that improves trust can nevertheless worsen equitable access. \\
Practical impacts & The framework can inform platform/ticketing design, while the open-access game supports SDG 4 learning. \\
Feedback received and revision made & The project expanded from reputation-only decisions to the interaction of reputation, platform rules, speculation, and income uncertainty. \\
\bottomrule
\end{tabularx}
\end{table}
